% !TeX program = lualatex
% !BIB program = biber
% Projeto UnDF-TeX. LPPL 1.3c ou posterior.
% Perfil experimental: conferir exigências de entrega com o curso.
\documentclass[perfil=artigo]{undftex}

\autor{Nome do Estudante}
\titulo{TCC em formato de artigo}
\subtitulo{modelo experimental do UnDF-TeX}
% \tituloartigoestrangeiro{Final paper in article format}
\curso{Nome do curso}
\tipotrabalho{Trabalho de Conclusão de Curso em formato de artigo}
\graupretendido{bacharel em Nome do curso}
\local{Brasília}
\ano{2026}
\orientador{Nome do Orientador}
\dataaprovacao{Brasília, DF, dia de mês de ano}
\limparmembrosbanca
\membrobanca{Nome do Orientador}{Prof. Dr.}{UnDF}
\membrobanca{Nome do Membro 1}{Profa. Dra.}{Instituição}
\membrobanca{Nome do Membro 2}{Prof. Dr.}{Instituição}
\addbibresource{obrigatorio/referencias.bib}

\undfresetpretextualcontents
% O mesmo resumo aparece na apresentação institucional e na abertura do miolo.
% A faixa de 150 a 250 palavras atende às duas recomendações de extensão.
% !TeX root = ../artigo.tex
% Texto demonstrativo: substitua por objetivo, método, resultados e conclusão reais.
\undfsetcontent{resumo}{%
Este modelo apresenta uma proposta de organização de trabalho de conclusão de
curso em formato de artigo científico ou tecnológico, independente de regras de
periódicos. O objetivo é oferecer uma base de edição com título, identificação da
autoria, resumo único, palavras-chave, seções contínuas e referências. A proposta
compartilha os mecanismos de citações, numeração e bibliografia do UnDF-TeX e
separa o conteúdo científico dos elementos de apresentação exigidos para entrega
institucional. O procedimento de verificação compreende compilação, conferência
das referências e exame visual dos documentos produzidos. Não são apresentados
resultados empíricos de pesquisa, e as escolhas de apresentação não pressupõem
aprovação por qualquer curso. O resultado esperado é um documento editável que
permita avaliar a comunicação concentrada da pesquisa e reutilizar o mesmo
conteúdo na entrega acadêmica. A adoção final depende da conferência das
orientações do curso, da banca e da biblioteca.}
\undfsetcontent{palavraschave}{trabalho de conclusão de curso; artigo científico; redação acadêmica}

\undfinputoptional{opcional/ficha-catalografica}

\begin{document}
  \undfvalidateonepagecontents
  \undfcover
  \pretextual
  \undffolhaderosto
  \undfversofolhaderostoopcional
  \undffolhadeaprovacao
  \begin{resumo}
    \undfusecontent{resumo}
    \palavraschave{\undfusecontent{palavraschave}}
  \end{resumo}
  \begin{resumoestrangeiro}
    \begin{otherlanguage}{english}
      \undfusecontent{abstract}
      \keywords{\undfusecontent{keywords}}
    \end{otherlanguage}
  \end{resumoestrangeiro}
  \undfsumario

  \undfinicioartigo
  % !TeX root = ../artigo.tex
% Substitua as orientações abaixo pelo conteúdo da pesquisa.
\section{Introdução}
\label{sec:artigo-introducao}
Delimite o tema, o problema, a justificativa e o objetivo do estudo. A redação
científica precisa relacionar a pergunta de pesquisa às evidências que serão
apresentadas, com fundamentação nas fontes efetivamente consultadas
\citeind{severino2017}. O recorte do artigo deve caber em uma exposição concisa
sem omitir o método nem as limitações do trabalho.

\section{Fundamentação e trabalhos relacionados}
Apresente os conceitos necessários e compare trabalhos relacionados ao problema.
Esta seção pode ser integrada à introdução quando isso favorecer a argumentação.
Evite transformar o artigo em uma sequência de resumos de autores.

\section{Método}
Descreva o delineamento, os materiais ou dados, os procedimentos de análise e os
cuidados éticos aplicáveis. Em artigo tecnológico, explique também o projeto e a
avaliação do artefato. Informe condições suficientes para compreender como as
evidências foram produzidas.

\subsection{Procedimentos de avaliação}
Defina os critérios, as métricas e as condições da avaliação antes de discutir os
resultados. Diferencie evidência observada de resultado esperado. O exemplo
matemático a seguir demonstra a numeração independente de capítulos:
\begin{equation}
  p = \frac{n_{\mathrm{observado}}}{n_{\mathrm{total}}}.
  \label{eq:artigo-proporcao}
\end{equation}
Defina cada grandeza e explique sua pertinência; remova esta demonstração quando
não houver cálculo correspondente no estudo.

\section{Resultados e discussão}
Apresente os achados, relacione-os ao objetivo e confronte-os com a literatura.
Resultados e discussão podem formar seções distintas. Discuta limitações e
ameaças à validade. Figuras e tabelas devem apoiar a análise e ter fonte indicada.

\section{Considerações finais}
Retome a resposta ao problema, a contribuição e os limites do estudo. Indique
desdobramentos fundamentados nos resultados, sem introduzir dados inéditos nesta
seção. Este arquivo contém orientações de redação e não resultados de uma pesquisa.

  \printundfreferences
  % Apêndices e anexos usam os comandos institucionais existentes.
\end{document}
